\documentclass[aps,pra,twocolumn,superscriptaddress,hyperref]{revtex4-2}
\usepackage{amsmath,amssymb,amsfonts,graphicx,booktabs}

\begin{document}

\title{Capacity-Normalized Bound on Non-Markovian Backflow from Finite Information Capacity}

\author{Huang Zhongchang}
\affiliation{Independent Researcher, Nanjing, China}

\begin{abstract}
Non-Markovian dynamics in open quantum systems---where information flows back from environment to system---is characterized by several complementary measures, from trace-distance revivals to entropic bounds. Here we derive a combinatorially simpler bound on the capacity-normalized backflow ratio: for a bipartite qubit system with a pure-dephasing pointer-basis interaction, $\mathcal{R} \equiv N_{\rm back}/C_F \leq q_S/q_E$ (for equal system and environment sizes $N_S=N_E$), where $C_F = N_E q_E$ is the forward transfer capacity, and $q_S$, $q_E$ are per-qubit populations of the undetermined state $|0\rangle$ before irreversible transfer begins. The bound follows from the pigeonhole principle---each qubit toggles $|0\rangle \to |1\rangle$ at most once---requires no spectral densities or coupling strengths, and is accessible via projective population measurements without full quantum state tomography. We illustrate with a six-qubit collision model and discuss conditions under which the bound is non-trivial. The bound is non-trivial when the system is more determined than the environment ($q_S < q_E$), the physically relevant regime for standard decoherence.
\end{abstract}

\maketitle

%%%%%%%%%%%%%%%%%%%%%%%%%%%%%%%%%%%%%%%%%%%%%%%%%%%%%%%%%%%%%%%%%%%%%%
\section{Introduction}
%%%%%%%%%%%%%%%%%%%%%%%%%%%%%%%%%%%%%%%%%%%%%%%%%%%%%%%%%%%%%%%%%%%%%%

Open quantum systems exhibit non-Markovian dynamics when information, having flowed from system to environment, subsequently returns. This backflow is detected by the Breuer-Laine-Piilo (BLP) measure~\cite{Breuer:2009}, which integrates positive excursions of the trace-distance derivative. The divisibility-based characterization of non-Markovianity was established by Rivas, Huelga, and Plenio (RHP)~\cite{Rivas:2010}, who proved that CP-divisibility of the dynamical map is a necessary and sufficient condition for Markovianity. An entropic upper bound on backflow was proved by Megier, Smirne, and Vacchini (MSV)~\cite{Megier:2021}, who showed that backflow is constrained by system-environment correlations quantified by the Holevo skew divergence. Buscemi et al.~\cite{Buscemi:2025} provide a process-matrix classification of causal versus noncausal revivals using quantum conditional mutual information, complementing the BLP--RHP--MSV landscape with a structural taxonomy.

This Letter supplies a combinatorially simpler bound. The MSV bound~\cite{Megier:2021} is general---it applies to any CPTP dynamics---but its evaluation requires quantum state tomography to reconstruct the joint system-environment state. The present bound is narrower in domain (it requires unidirectional Lindblad dynamics) but operationally simpler: it depends only on pointer-basis population averages, accessible via projective measurements without tomography.

We derive the bound from a single physical assumption: each qubit carries at most one bit of distinguishable information in the einselected pointer basis. Forward transfer is described by a Lindblad jump operator that toggles $|0\rangle \to |1\rangle$ without a reverse channel. Under these conditions, the capacity-normalized backflow ratio $\mathcal{R} \equiv N_{\rm back}/C_F$ satisfies $\mathcal{R} \leq q_S/q_E$ (for equal system-environment sizes). The proof requires only counting.

%%%%%%%%%%%%%%%%%%%%%%%%%%%%%%%%%%%%%%%%%%%%%%%%%%%%%%%%%%%%%%%%%%%%%%
\section{Model}
\label{sec:model}
%%%%%%%%%%%%%%%%%%%%%%%%%%%%%%%%%%%%%%%%%%%%%%%%%%%%%%%%%%%%%%%%%%%%%%

\subsection{Einselection of the pointer basis}

A dominant pure-dephasing interaction
\begin{equation}
H_{\rm int} = \sum_{i} \sigma_z^{(i)} \otimes B_i,
\label{eq:hint}
\end{equation}
with $\sigma_z^{(i)} = |0\rangle\langle 0| - |1\rangle\langle 1|$, einselects the pointer basis $\{|0\rangle,|1\rangle\}$ via environment-induced superselection~\cite{Zurek:2003}. This interaction preserves pointer-basis populations; its sole role is basis selection. Population dynamics are supplied by the Lindblad channel below.

\subsection{Binary information capacity}

In the einselected pointer basis, each qubit carries at most one bit of distinguishable classical information---the finite-capacity hypothesis. The state $|0\rangle$ is undetermined: no causal event has fixed the qubit's pointer value. The state $|1\rangle$ is determined: a causal event has fixed it. A binary pointer basis accessed via projective measurement thus yields at most one classical bit~\cite{Zurek:2003}. This is a consequence of einselection, not an additional postulate: the pointer basis is the only basis in which classical information can be stably encoded, and it supports exactly two distinguishable states.

\subsection{Forward transfer}

Forward transfer---a determined qubit influencing an undetermined neighbor---is described by the Lindblad jump operator~\cite{Lindblad:1976,GKS:1976}
\begin{equation}
L_{S\to E} = |1\rangle\langle 1|_S \otimes |1\rangle\langle 0|_E,
\label{eq:jump}
\end{equation}
acting on a specific $S$--$E$ pair. The jump toggles the target to $|1\rangle$ when the source is $|1\rangle$ and the target $|0\rangle$; otherwise it is null. The source qubit is unchanged.

Equation~\eqref{eq:jump} creates $|1\rangle$ without a reverse channel. The absence of a $|0\rangle\langle 1|$ term reflects the energy asymmetry $\Delta E \gg k_B T$, under which spontaneous de-determination is exponentially suppressed. Neither $H_{\rm int}$ nor $L$ contains a $|0\rangle\langle 1|$ channel, so the GKSL dynamics preserves the directedness of $|0\rangle \to |1\rangle$. When $|1\rangle \to |0\rangle$ relaxation is non-negligible, $T_1$ decay creates apparent undetermined $S$-qubits that were previously determined, inflating the effective $q_S$; the bound then requires a $T_1$ correction.

%%%%%%%%%%%%%%%%%%%%%%%%%%%%%%%%%%%%%%%%%%%%%%%%%%%%%%%%%%%%%%%%%%%%%%
\section{Bound on Capacity-Normalized Backflow}
\label{sec:theorem}
%%%%%%%%%%%%%%%%%%%%%%%%%%%%%%%%%%%%%%%%%%%%%%%%%%%%%%%%%%%%%%%%%%%%%%

\subsection{Definitions}

Partition the qubits into a system $S$ ($N_S$ qubits) and an environment $E$ ($N_E$ qubits). Define the pre-transfer pointer-basis populations:
\begin{equation}
q_S = \frac{1}{N_S}\sum_{i\in S} \langle 0|\rho_i|0\rangle, \qquad
q_E = \frac{1}{N_E}\sum_{i\in E} \langle 0|\rho_i|0\rangle.
\end{equation}

The \textbf{forward transfer capacity} is the number of $E$-qubits initially available to receive determination:
\begin{equation}
C_F \equiv N_E q_E.
\end{equation}

A \textbf{backflow event} is an $E \to S$ determination transfer: an $E$-qubit in $|1\rangle$ toggles an undetermined $S$-qubit from $|0\rangle$ to $|1\rangle$ via $L_{E\to S}$. Let $N_{\rm back}$ be the total count of such events. The \textbf{capacity-normalized backflow ratio} is:
\begin{equation}
\mathcal{R} \equiv \frac{N_{\rm back}}{C_F}.
\end{equation}

$\mathcal{R}$ quantifies how efficiently the environment converts its forward-transfer capacity into return flow: $\mathcal{R}=1$ means every unit of forward capacity is matched by a backflow event, saturating the bound. $\mathcal{R}$ is not the fraction of forward-transfer events that reverse (which would be $N_{\rm back}/N_F$, where $N_F$ is the actual number of forward events). Since $N_F \leq C_F$, we have $N_{\rm back}/N_F \geq \mathcal{R}$: the actual reflux fraction is at least as large as $\mathcal{R}$. The bound on $\mathcal{R}$ therefore provides a lower bound on the actual reflux fraction when $N_F$ is known, though a tight upper bound on $N_{\rm back}/N_F$ would require additional dynamical assumptions beyond the counting argument.

\subsection{Theorem and proof}

\begin{quote}
\textbf{Theorem (Capacity-Normalized Backflow Bound).}
$\mathcal{R} \leq N_S q_S / (N_E q_E)$.
For $N_S = N_E$: $\mathcal{R} \leq q_S/q_E$.
\end{quote}

\textbf{Proof.}
Each qubit toggles $|0\rangle \to |1\rangle$ at most once in the pointer basis, since $L$ contains no reverse channel and $H_{\rm int}$ preserves pointer populations. An $S$-qubit that toggles during backflow must be initially undetermined; there are at most $N_S q_S$ such qubits on average. In any single experimental run the integer count is $\lfloor N_S q_S \rfloor$ or $\lceil N_S q_S \rceil$; the bound holds in expectation and with high probability for $N_S \gg 1$. Hence $N_{\rm back} \leq N_S q_S$. By definition $\mathcal{R} \equiv N_{\rm back}/C_F = N_{\rm back}/(N_E q_E)$. Therefore $\mathcal{R} \leq N_S q_S/(N_E q_E)$. For $N_S = N_E$, this reduces to $q_S/q_E$. $\square$

The inequality is the pigeonhole principle applied to qubit toggle events: exceeding $N_S q_S$ backflow events would force an $S$-qubit to toggle twice, which the unidirectional dynamics forbids. The bound applies to ensemble-averaged counts; in any single run, integer fluctuations around $N_S q_S$ leave the inequality unchanged in expectation. The bound depends on no spectral densities, coupling strengths, or noise models; it requires only the finite-capacity hypothesis and the directedness of $L$.

\subsection{Domain and non-triviality}

The bound requires $q_E > 0$ (otherwise $C_F = 0$ and $\mathcal{R}$ is undefined). It is non-trivial ($\mathcal{R} < 1$) when $q_S < q_E$: the system must be more determined than the environment. When $q_S \geq q_E$, the bound exceeds unity and provides no constraint.

The general form $\mathcal{R} \leq N_S q_S/(N_E q_E)$ is tighter for $N_S \ll N_E$: a small system coupled to a large environment has $\mathcal{R} \to 0$, consistent with the effective irreversibility of standard decoherence~\cite{Breuer:2016}.

\subsection{$T_1$ correction}

When $|1\rangle \to |0\rangle$ relaxation is non-negligible (finite $T_1$), the bound acquires a correction. An $S$-qubit that decays $|1\rangle \to |0\rangle$ before measurement is misidentified as undetermined, inflating the apparent $q_S$. The corrected bound is $\mathcal{R} \leq q_S/q_E + \delta$, where $\delta \sim \tau_{\rm prot}/T_1$ and $\tau_{\rm prot}$ is the total protocol duration. A discussion of gate-level considerations for experimental implementation is provided in the Supplemental Material~\cite{SM}.

%%%%%%%%%%%%%%%%%%%%%%%%%%%%%%%%%%%%%%%%%%%%%%%%%%%%%%%%%%%%%%%%%%%%%%
\section{Relation to Prior Work}
%%%%%%%%%%%%%%%%%%%%%%%%%%%%%%%%%%%%%%%%%%%%%%%%%%%%%%%%%%%%%%%%%%%%%%

The MSV entropic bound~\cite{Megier:2021} constrains backflow via the Holevo skew divergence $S_\mu(\varrho_{SE}, \varrho_S \otimes \varrho_E)$ and applies to any CPTP dynamics. Its evaluation requires full state tomography. The present bound $\mathcal{R} \leq q_S/q_E$ is narrower in domain (requires unidirectional $L$) but requires only projective population measurements. The two bounds are complementary: MSV for dynamical generality, ours for operational simplicity.

The BLP measure~\cite{Breuer:2009, Breuer:2016} detects non-Markovianity by integrating $\dot{D}(t) > 0$. Rivas, Huelga, and Plenio~\cite{Rivas:2010} established the CP-divisibility criterion, proving equivalence between Markovianity and the CP character of intermediate dynamical maps. The connection between trace-distance revivals and system-environment correlations was established by Laine, Piilo, and Breuer~\cite{Laine:2010}. For a comprehensive review of these and other measures, see Breuer~\cite{Breuer:2012}. Buscemi et al.~\cite{Buscemi:2025} classify causal versus noncausal revivals within the process-matrix framework, addressing a question orthogonal to ours: whether a revival is causal or not, rather than how large backflow can be. The present work adds a combinatorially derived, operationally accessible bound on the capacity-normalized backflow ratio.

%%%%%%%%%%%%%%%%%%%%%%%%%%%%%%%%%%%%%%%%%%%%%%%%%%%%%%%%%%%%%%%%%%%%%%
\section{Illustration: Six-Qubit Collision Model}
%%%%%%%%%%%%%%%%%%%%%%%%%%%%%%%%%%%%%%%%%%%%%%%%%%%%%%%%%%%%%%%%%%%%%%

Consider a collision model with $N_S = N_E = 3$ qubits and all-to-all $S$--$E$ connectivity. Figure~\ref{fig:bound} shows the bound $\mathcal{R} \leq q_S/q_E$ as a function of $q_E$ for several values of $q_S$, with the trivial boundary ($\mathcal{R}=1$) marked. The non-trivial regime $q_S < q_E$ is shaded.

Table~\ref{tab:examples} provides representative $(q_S, q_E)$ pairs with their physical interpretation. The bound tightens as $q_S/q_E$ decreases, i.e., as the system becomes more determined relative to the environment. The bound is saturated when all undetermined $E$-qubits receive forward transfer and all undetermined $S$-qubits receive backflow. In the collision model with uniform all-to-all coupling, this saturation is approachable when the number of applied gates is large compared to the number of undetermined qubits.

\begin{table}[t]
\centering
\caption{Representative $\mathcal{R}$ bounds for $N_S = N_E = 3$. The non-trivial regime is $q_S < q_E$ (system more determined than environment).}
\label{tab:examples}
\begin{tabular}{cccl}
\toprule
$q_S$ & $q_E$ & $\mathcal{R} \leq q_S/q_E$ & Physical interpretation \\
\midrule
0.1 & 0.9 & 0.11 & Nearly-determined system, fresh environment \\
0.2 & 0.8 & 0.25 & Moderately determined system \\
0.3 & 0.6 & 0.50 & Slightly more determined system \\
0.4 & 0.5 & 0.80 & Marginally non-trivial \\
0.5 & 0.5 & 1.00 & Trivial (equal determination) \\
0.7 & 0.3 & 2.33 & Trivial ($q_S > q_E$, no constraint) \\
0.0 & 0.5 & 0.00 & Fully determined system \\
\bottomrule
\end{tabular}
\end{table}

\begin{figure}[t]
\centering
\includegraphics[width=\columnwidth]{fig1_bound.pdf}
\caption{Capacity-normalized backflow bound $\mathcal{R} \leq q_S/q_E$ as a function of $q_E$ for representative values of $q_S$. The dashed line marks the trivial boundary $\mathcal{R}=1$. The shaded region $q_S < q_E$ is the non-trivial regime where the bound constrains backflow.}
\label{fig:bound}
\end{figure}

%%%%%%%%%%%%%%%%%%%%%%%%%%%%%%%%%%%%%%%%%%%%%%%%%%%%%%%%%%%%%%%%%%%%%%
\section{Discussion}
%%%%%%%%%%%%%%%%%%%%%%%%%%%%%%%%%%%%%%%%%%%%%%%%%%%%%%%%%%%%%%%%%%%%%%

The counting argument is classical---binary distinguishability and the pigeonhole principle. Its applicability to quantum systems rests on the quantum-mechanical origin of the pointer basis via einselection~\cite{Zurek:2003} and the directedness of determination from the GKSL form of $L$~\cite{Lindblad:1976,GKS:1976}. Quantum mechanics supplies the physical conditions under which a combinatorial inequality constrains a real system.

The bound $\mathcal{R} \leq q_S/q_E$ is a necessary condition for any physical process consistent with the model assumptions. A violation---$\mathcal{R}$ exceeding $q_S/q_E$---would indicate either (a) a qubit toggled $|0\rangle \to |1\rangle$ more than once (gate error or de-determination), (b) $T_1$ decay inflating apparent $q_S$ beyond the correction $\delta$, or (c) a failure of the finite-capacity hypothesis. Distinguishing these failure modes requires independent gate characterization, discussed in the Supplemental Material~\cite{SM}.

The bound's principal limitation is that it constrains $\mathcal{R} \equiv N_{\rm back}/C_F$, not the actual reflux fraction $N_{\rm back}/N_F$. The actual reflux fraction can exceed $\mathcal{R}$ by a factor of $C_F/N_F$, which may be large when few forward-transfer events occur. The capacity-normalized quantity is chosen because it admits a strict combinatorial bound; an equally strict bound on $N_{\rm back}/N_F$ would require additional dynamical assumptions linking $N_F$ to $C_F$.

Beyond the present model, extending the bound to bidirectional dynamics---where $|0\rangle\langle 1|$ channels are present---would require a generalized counting argument that tracks net toggle events rather than absolute counts. In that setting each qubit can toggle in both directions, and the simple bound $N_{\rm back} \leq N_S q_S$ no longer holds. A net-event formulation, possibly involving a detailed-balance relation between forward and reverse rates, is a natural direction for future work.

%%%%%%%%%%%%%%%%%%%%%%%%%%%%%%%%%%%%%%%%%%%%%%%%%%%%%%%%%%%%%%%%%%%%%%
\section{Conclusion}
%%%%%%%%%%%%%%%%%%%%%%%%%%%%%%%%%%%%%%%%%%%%%%%%%%%%%%%%%%%%%%%%%%%%%%

Under the finite-capacity hypothesis (each qubit $\leq 1$ bit in its pointer basis), the capacity-normalized backflow ratio satisfies $\mathcal{R} \equiv N_{\rm back}/C_F \leq N_S q_S/(N_E q_E)$. The bound follows rigorously from the model assumptions, is operationally accessible via projective population measurements, and is complementary to the entropic bound of Megier et al.~\cite{Megier:2021}. It provides a simple quantitative constraint on backflow in open quantum systems with unidirectional Lindblad dynamics, linking the finite information capacity of the physical substrate to observable consequences in the pointer basis.

\begin{thebibliography}{14}

\bibitem{Breuer:2009}
H.-P.~Breuer, E.-M.~Laine, and J.~Piilo,
Measure for the degree of non-Markovian behavior,
Phys.\ Rev.\ Lett.\ \textbf{103}, 210401 (2009).

\bibitem{Breuer:2016}
H.-P.~Breuer, E.-M.~Laine, J.~Piilo, and B.~Vacchini,
Colloquium: Non-Markovian dynamics in open quantum systems,
Rev.\ Mod.\ Phys.\ \textbf{88}, 021002 (2016).

\bibitem{Rivas:2010}
\'A.~Rivas, S.~F.~Huelga, and M.~B.~Plenio,
Entanglement and non-Markovianity of quantum evolutions,
Phys.\ Rev.\ Lett.\ \textbf{105}, 050403 (2010).

\bibitem{Laine:2010}
E.-M.~Laine, J.~Piilo, and H.-P.~Breuer,
Measure for the non-Markovianity of quantum processes,
Europhys.\ Lett.\ \textbf{92}, 60010 (2010).

\bibitem{Breuer:2012}
H.-P.~Breuer,
Foundations and measures of quantum non-Markovianity,
J.\ Phys.\ B \textbf{45}, 154001 (2012).

\bibitem{Buscemi:2025}
F.~Buscemi, R.~Gangwar, K.~Goswami, H.~Badhani, T.~Pandit, B.~Mohan, S.~Das, and M.~N.~Bera,
Causal and noncausal revivals of information,
PRX Quantum \textbf{6}, 020316 (2025).

\bibitem{Megier:2021}
N.~Megier, A.~Smirne, and B.~Vacchini,
Entropic bounds on information backflow,
Phys.\ Rev.\ Lett.\ \textbf{127}, 030401 (2021).

\bibitem{Zurek:2003}
W.~H.~Zurek,
Decoherence, einselection, and the quantum origins of the classical,
Rev.\ Mod.\ Phys.\ \textbf{75}, 715 (2003).

\bibitem{Lindblad:1976}
G.~Lindblad,
On the generators of quantum dynamical semigroups,
Commun.\ Math.\ Phys.\ \textbf{48}, 119 (1976).

\bibitem{GKS:1976}
V.~Gorini, A.~Kossakowski, and E.~C.~G.~Sudarshan,
Completely positive dynamical semigroups of $N$-level systems,
J.\ Math.\ Phys.\ \textbf{17}, 821 (1976).

\bibitem{SM}
See Supplemental Material for the complete proof with domain conditions, the physical origin of $L$, comparison to the MSV bound, a numerical illustration, and a quantitative discussion of gate-level considerations for experimental implementation.

\end{thebibliography}

\acknowledgments
AI language tools were used for editing assistance.

\end{document}
